\addtocontents{toc}{\protect\setcounter{tocdepth}{0}}
\chapter[Anexo: Manual de Usuario]{ANEXO: Manual de Usuario}

Este manual explica el uso de OnCuidar para el cuidador principal. Cada sección
describe una tarea de principio a fin.

\section{Crear una cuenta e ingresar}

Al abrir la aplicación por primera vez aparece la pantalla de bienvenida.
Presione \emph{Crear cuenta} y complete sus datos, los del paciente, los del
centro de salud y los del contacto de emergencia. Al terminar, la aplicación lo
lleva al panel principal. En los usos siguientes la sesión se mantiene abierta.
Si cierra sesión, ingrese con su correo y contraseña. Si olvidó la contraseña,
presione \emph{¿Olvidaste tu contraseña?} y siga el enlace que llegará a su
correo.

\TBDFigura{capturas de la bienvenida, la creación de cuenta y el inicio de sesión}

\section{Confirmar su correo}

Después de crear la cuenta, OnCuidar envía un mensaje con un enlace de
verificación al correo registrado. Abra el mensaje y presione el enlace; se
abrirá una página que indica \emph{Correo confirmado}. La confirmación no es
obligatoria para usar la aplicación, pero asegura que podrá recuperar el acceso.
Si el mensaje no llega, revise la carpeta de correo no deseado y, en
\emph{Perfil}, solicite que se reenvíe el enlace (se puede pedir uno por minuto).
El enlace sirve una sola vez y vence a la hora de enviado.

Para cambiar su correo principal o agregar un correo de respaldo, vaya a
\emph{Perfil}, ingrese el correo nuevo y su contraseña actual. El cambio queda
pendiente y solo se aplica cuando presiona el enlace enviado al correo nuevo;
mientras tanto sigue vigente el correo anterior.

\TBDFigura{captura de la insignia de correo verificado en el perfil}

\section{Recuperar el acceso}

Si olvidó su contraseña, en la pantalla de inicio de sesión presione
\emph{¿Olvidaste tu contraseña?} e ingrese su correo; también puede usar el
correo de respaldo si lo confirmó. La aplicación siempre responde que el mensaje
fue enviado, aunque el correo no exista, para no revelar qué cuentas hay. Abra
el enlace que recibe, escriba la contraseña nueva en la página que se muestra y
vuelva a la aplicación para ingresar. El enlace sirve una sola vez y vence a la
hora. Al iniciar sesión, incluso en un teléfono nuevo, la aplicación recupera su
clave de cifrado desde el servidor y sus datos vuelven a estar disponibles.

\TBDFigura{capturas de la recuperación de la contraseña}

\section{Registrar la valoración del paciente}

Presione \emph{Registro} en la barra inferior. Ingrese los signos vitales que
tenga disponibles, mueva el indicador de cada síntoma según su intensidad de 0 a
10 (0 sin síntoma, 1 a 3 leve, 4 a 6 moderado, 7 a 9 severo y 10 insoportable)
y agregue observaciones si lo desea. Al guardar, la aplicación muestra el
nivel de alerta.

\begin{itemize}
    \item \emph{Verde (normal)}: continúe con la vigilancia habitual.
    \item \emph{Amarillo (alerta)}: se recomienda consultar al equipo de salud.
    \item \emph{Rojo (crítico)}: contacte al equipo de salud de inmediato.
\end{itemize}

La aplicación orienta, pero no reemplaza la evaluación de un profesional. Ante
cualquier duda, comuníquese con el equipo tratante.

\TBDFigura{captura del registro clínico y del resultado de la alerta}

\section{Revisar y exportar el historial}

Desde el panel principal abra el \emph{Historial}. Los registros aparecen del más
reciente al más antiguo y puede filtrarlos por fecha o nivel. Para llevar el
historial a un control, presione exportar, elija PDF o Excel, indique el rango de
fechas y luego abra o comparta el archivo.

\TBDFigura{captura del historial y de la exportación}

\section{Resolver dudas}

En \emph{Chat} escriba su duda con sus propias palabras o elija una de las
preguntas sugeridas. El asistente responde con información revisada por el
equipo de Enfermería. Si no encuentra una respuesta, le sugerirá los temas
disponibles. También puede buscar en las \emph{Preguntas frecuentes} por
palabra clave.

\TBDFigura{captura del chat y de las preguntas frecuentes}

\section{Aprender con la biblioteca}

En \emph{Aprende} encontrará videos, guías, infografías y documentos. Filtre por
categoría o busque por título, y marque como favoritos los materiales que quiera
tener a mano. El contenido se descarga automáticamente, por lo que puede
consultarlo sin conexión.

\TBDFigura{captura de la biblioteca educativa}

\section{Programar recordatorios}

Desde el panel principal abra \emph{Recordatorios} y cree uno con su título,
tipo, fecha, hora y repetición. El teléfono le avisará a la hora indicada,
aunque la aplicación esté cerrada. Para eliminar un recordatorio, use el menú de
la tarjeta.

\TBDFigura{captura de los recordatorios}

\section{Perfil y contacto con el centro de salud}

En \emph{Perfil} puede actualizar sus datos, agregar pacientes y cambiar el
paciente activo. Desde la tarjeta del paciente puede llamar al centro de salud o
abrir su dirección en el mapa.

\TBDFigura{captura del perfil}
